\section{Generically finite morphisms}
\label{section-generically-finite}

\noindent
In this section we revisit the notion of a generically finite
morphism of schemes as studied in
Morphisms, Section \ref{morphisms-section-generically-finite}.

\begin{lemma}
\label{lemma-quasi-finite-in-codim-1}
Let $f : X \to Y$ be locally of finite type. Let $y \in Y$ be a point
such that $\mathcal{O}_{Y, y}$ is Noetherian of dimension $\leq 1$.
Assume in addition one of the following conditions is satisfied
\begin{enumerate}
\item for every generic point $\eta$ of an irreducible component
of $X$ the field extension $\kappa(\eta)/\kappa(f(\eta))$
is finite (or algebraic),
\item for every generic point $\eta$ of an irreducible component
of $X$ such that $f(\eta) \leadsto y$ the field extension
$\kappa(\eta)/\kappa(f(\eta))$ is finite (or algebraic),
\item $f$ is quasi-finite at every generic point of an
irreducible component of $X$,
\item $Y$ is locally Noetherian and $f$
is quasi-finite at a dense set of points of $X$,
\item add more here.
\end{enumerate}
Then $f$ is quasi-finite at every point of $X$ lying over $y$.
\end{lemma}

\begin{proof}
Condition (4) implies $X$ is locally Noetherian
(Morphisms, Lemma \ref{morphisms-lemma-finite-type-noetherian}).
The set of points at which morphism is quasi-finite is open
(Morphisms, Lemma \ref{morphisms-lemma-quasi-finite-points-open}).
A dense open of a locally Noetherian scheme contains all generic
point of irreducible components, hence (4) implies (3).
Condition (3) implies condition
(1) by Morphisms, Lemma \ref{morphisms-lemma-residue-field-quasi-finite}.
Condition (1) implies condition (2).
Thus it suffices to prove the lemma in case (2) holds.

\medskip\noindent
Assume (2) holds. Recall that $\Spec(\mathcal{O}_{Y, y})$
is the set of points of $Y$ specializing to $y$, see
Schemes, Lemma \ref{schemes-lemma-specialize-points}.
Combined with
Morphisms, Lemma \ref{morphisms-lemma-base-change-quasi-finite}
this shows we may replace $Y$ by $\Spec(\mathcal{O}_{Y, y})$.
Thus we may assume $Y = \Spec(B)$ where $B$ is a Noetherian
local ring of dimension $\leq 1$ and $y$ is the closed point.

\medskip\noindent
Let $X = \bigcup X_i$ be the irreducible components of $X$ viewed
as reduced closed subschemes. If we can show each fibre $X_{i, y}$
is a discrete space, then $X_y = \bigcup X_{i, y}$ is discrete as
well and we conclude that $X \to Y$ is quasi-finite at all points
of $X_y$ by Morphisms, Lemma
\ref{morphisms-lemma-quasi-finite-at-point-characterize}.
Thus we may assume $X$ is an integral scheme.

\medskip\noindent
If $X \to Y$ maps the generic point $\eta$ of $X$ to $y$, then $X$
is the spectrum of a finite extension of $\kappa(y)$ and the
result is true. Assume that $X$ maps $\eta$ to a point corresponding
to a minimal prime $\mathfrak q$ of $B$ different from $\mathfrak m_B$.
We obtain a factorization $X \to \Spec(B/\mathfrak q) \to \Spec(B)$.
Let $x \in X$ be a point lying over $y$.
By the dimension formula
(Morphisms, Lemma \ref{morphisms-lemma-dimension-formula})
we have
$$
\dim(\mathcal{O}_{X, x}) \leq \dim(B/\mathfrak q) +
\text{trdeg}_{\kappa(\mathfrak q)}(R(X)) - \text{trdeg}_{\kappa(y)} \kappa(x)
$$
We know that $\dim(B/\mathfrak q) = 1$, that the generic point of $X$
is not equal to $x$ and specializes to $x$ and that $R(X)$ is algebraic
over $\kappa(\mathfrak q)$. Thus we get
$$
1 \leq 1 - \text{trdeg}_{\kappa(y)} \kappa(x)
$$
Hence every point $x$ of $X_y$ is closed in $X_y$ by
Morphisms, Lemma
\ref{morphisms-lemma-algebraic-residue-field-extension-closed-point-fibre}
and hence $X \to Y$ is quasi-finite at every point $x$ of $X_y$ by
Morphisms, Lemma \ref{morphisms-lemma-quasi-finite-at-point-characterize}
(which also implies that $X_y$ is a discrete topological space).
\end{proof}

\begin{lemma}
\label{lemma-finite-in-codim-1}
Let $f : X \to Y$ be a proper morphism. Let $y \in Y$ be a point
such that $\mathcal{O}_{Y, y}$ is Noetherian of dimension $\leq 1$.
Assume in addition one of the following conditions is satisfied
\begin{enumerate}
\item for every generic point $\eta$ of an irreducible component
of $X$ the field extension $\kappa(\eta)/\kappa(f(\eta))$
is finite (or algebraic),
\item for every generic point $\eta$ of an irreducible component
of $X$ such that $f(\eta) \leadsto y$ the field extension
$\kappa(\eta)/\kappa(f(\eta))$ is finite (or algebraic),
\item $f$ is quasi-finite at every generic point of $X$,
\item $Y$ is locally Noetherian and $f$
is quasi-finite at a dense set of points of $X$,
\item add more here.
\end{enumerate}
Then there exists an open neighbourhood $V \subset Y$ of $y$ such that
$f^{-1}(V) \to V$ is finite.
\end{lemma}

\begin{proof}
By Lemma \ref{lemma-quasi-finite-in-codim-1} the morphism $f$ is
quasi-finite at every point of the fibre $X_y$. Hence
$X_y$ is a discrete topological space
(Morphisms, Lemma \ref{morphisms-lemma-quasi-finite-at-point-characterize}).
As $f$ is proper the fibre $X_y$ is quasi-compact, i.e., finite.
Thus we can apply Cohomology of Schemes, Lemma
\ref{coherent-lemma-proper-finite-fibre-finite-in-neighbourhood}
to conclude.
\end{proof}

\begin{lemma}
\label{lemma-modification-normal-iso-over-codimension-1}
Let $X$ be a Noetherian scheme. Let $f : Y \to X$ be a birational proper
morphism of schemes with $Y$ reduced. Let $U \subset X$ be the
maximal open over which $f$ is an isomorphism. Then $U$ contains
\begin{enumerate}
\item every point of codimension $0$ in $X$,
\item every $x \in X$ of codimension $1$ on $X$ such that
$\mathcal{O}_{X, x}$ is a discrete valuation ring,
\item every $x \in X$ such that the fibre of $Y \to X$ over $x$ is
finite and such that $\mathcal{O}_{X, x}$ is normal, and
\item every $x \in X$ such that $f$ is quasi-finite at some
$y \in Y$ lying over $x$ and $\mathcal{O}_{X, x}$ is normal.
\end{enumerate}
\end{lemma}

\begin{proof}
Part (1) follows from Morphisms, Lemma
\ref{morphisms-lemma-birational-isomorphism-over-dense-open}.
Part (2) follows from part (3) and Lemma \ref{lemma-finite-in-codim-1}
(and the fact that finite morphisms have finite fibres).

\medskip\noindent
Part (3) follows from part (4) and
Morphisms, Lemma \ref{morphisms-lemma-finite-fibre}
but we will also give a direct proof.
Let $x \in X$ be as in (3). By
Cohomology of Schemes, Lemma
\ref{coherent-lemma-proper-finite-fibre-finite-in-neighbourhood}
we may assume $f$ is finite. We may assume $X$ affine.
This reduces us to
the case of a finite birational morphism of Noetherian affine schemes
$Y \to X$ and $x \in X$ such that $\mathcal{O}_{X, x}$ is a
normal domain. Since $\mathcal{O}_{X, x}$ is a domain and $X$
is Noetherian, we may replace $X$ by an affine open of $x$ which
is integral. Then, since $Y \to X$ is birational and $Y$ is reduced
we see that $Y$ is integral. Writing $X = \Spec(A)$ and $Y = \Spec(B)$
we see that $A \subset B$ is a finite inclusion of domains having the same
field of fractions. If $\mathfrak p \subset A$ is the prime corresponding
to $x$, then $A_\mathfrak p$ being normal implies that
$A_\mathfrak p \subset B_\mathfrak p$ is an equality.
Since $B$ is a finite $A$-module, we see there exists an
$a \in A$, $a \not \in \mathfrak p$ such that $A_a \to B_a$
is an isomorphism.

\medskip\noindent
Let $x \in X$ and $y \in Y$ be as in (4). After replacing $X$
by an affine open neighbourhood we may assume $X = \Spec(A)$
and $A \subset \mathcal{O}_{X, x}$, see
Properties, Lemma \ref{properties-lemma-ring-affine-open-injective-local-ring}.
Then $A$ is a domain and hence $X$ is integral.
Since $f$ is birational and $Y$ is reduced
it follows that $Y$ is integral too. Consider the ring map
$\mathcal{O}_{X, x} \to \mathcal{O}_{Y, y}$. This is a ring map
which is essentially of finite type, the residue field extension
is finite, and $\dim(\mathcal{O}_{Y, y}/\mathfrak m_x\mathcal{O}_{Y, y}) = 0$
(to see this trace through the definitions of quasi-finite
maps in
Morphisms, Definition \ref{morphisms-definition-quasi-finite} and
Algebra, Definition \ref{algebra-definition-quasi-finite}). By
Algebra, Lemma \ref{algebra-lemma-essentially-finite-type-fibre-dim-zero}
$\mathcal{O}_{Y, y}$ is the localization of a finite
$\mathcal{O}_{X, x}$-algebra $B$. Of course we may replace
$B$ by the image of $B$ in $\mathcal{O}_{Y, y}$ and assume
that $B$ is a domain with the same fraction field as $\mathcal{O}_{Y, y}$.
Then $\mathcal{O}_{X, x} \subset B$
have the same fraction field as $f$ is birational. Since
$\mathcal{O}_{X, x}$ is normal, we conclude that
$\mathcal{O}_{X, x} = B$ (because finite implies integral),
in particular, we see that $\mathcal{O}_{X, x} = \mathcal{O}_{Y, y}$. By
Morphisms, Lemma \ref{morphisms-lemma-morphism-defined-local-ring}
after shrinking $X$ we may assume there is a section
$X \to Y$ of $f$ mapping $x$ to $y$ and inducing the given
isomorphism on local rings. Since $X \to Y$ is closed
(by Schemes, Lemma \ref{schemes-lemma-section-immersion})
necessarily maps the generic point of $X$ to the generic point of $Y$
it follows that the image of $X \to Y$ is $Y$.
Then $Y = X$ and we've proved what we wanted to show.
\end{proof}






